\section{Sample calculation at one-loop order}\label{sec:sample}

For illustration, the divergent parts of the one-loop diagrams
that contribute to the gauge-field two-point functions are
worked out in this section. The parameter and field
renormalization required at flow time zero
is then seen to cancel the singularities at all flow times.

\begin{figure}[t]
\centerline{\includegraphics[width=9.6cm]{images/diagram7}}
\caption{One-loop diagrams contributing to the
$AA$ ($1-3$), $LA$ ($4$), $LB$ ($5,6$) and
$LL$ ($7$) vertex functions.
The time-momentum-index combinations at the
external legs on the left and right
are $(t,p,\mu,a)$ and $(s,q,\nu,b)$, respectively
(the flow times are absent in the case of the $A$ legs).
}
\label{fig:oneloop}
\end{figure}

\subsection{Computation of self-energy diagrams}\label{ssec:selfenergy}

There are only few self-energy diagrams at one-loop order (see fig.~\ref{fig:oneloop}).
The diagrams $1-3$ coincide with the ones contributing to the
two-point function of the gauge field at flow time zero.
We therefore merely quote the known result
\begin{equation}
   \left.\Gamma^{(1)}_{AA}(p)_{\mu\nu}^{ab}
   \right|_{\rm pole}
   =\delta^{ab}(\delta_{\mu\nu}p^2-p_{\mu}p_{\nu})
   {N\over16\pi^2\eps}\left({13\over6}-{1\over 2\lambda_0}\right)
   \label{eq:5.1}
\end{equation}
for the sum of their singular parts.

In the case of diagram $4$,
the Feynman integrand is a linear combination
of terms of the form
\begin{equation}
   \delta^{ab}
   {Q(k,p)_{\mu\nu}\over \bigl\{k^2(k+p)^2\bigr\}^2}
   \rme^{-t\{uk^2+v(k+p)^2\}},
   \qquad u,v\in\{1,\alpha_0\},
   \label{eq:5.2}
\end{equation}
where $k$ is the loop momentum and $Q(k,p)_{\mu\nu}$
a homogeneous polynomial in $k$ and $p$ of degree $6$
(the momentum-conservation factor $(2\pi)^D\delta(p+q)$
is now always suppressed). This diagram is finite
at all $t>0$ and in any dimension $D$,
but eventually it will be multiplied by
the external propagators and must then be integrated
over $t$ from $0$ to infinity. The behaviour
of the integral small $t$ thus matters and can give rise
to singularities at $D=4$.

If $f(t)$ is any smooth test function (an external propagator, for
example), the time integral
\begin{equation}
   I(k,p)=\int_{0}^{\infty}\rmd t\, f(t)\rme^{-t\{uk^2+v(k+p)^2\}}
\label{eq:5.3}
\end{equation}
can be worked out in an asymptotic series at large $k$,
the leading term being
\begin{equation}
   I(k,p)={f(0)\over(u+v)k^2}+\rmO(k^{-3}).
\label{eq:5.4}
\end{equation}
The integration over the flow time $t$ thus improves the degree of
divergence of the momentum integral. In particular, the first term
in the decomposition
\begin{equation}
   I(k,p)=\left\{I(k,p)-{f(0)\over(u+v)(k+r)^2}\right\}+
   {f(0)\over(u+v)(k+r)^2}
   \label{eq:5.5}
\end{equation}
(where $r$ is an arbitrary external momentum) makes the integral
convergent at $D=4$.
One is then left with a logarithmic divergence and
therefore a pole singularity equal to
the one of the ordinary Feynman integral
\begin{equation}
   \delta^{ab}\delta(t){1\over u+v}\int_k
   {Q(k,0)_{\mu\nu}\over(k^2)^4(k+r)^2},
   \label{eq:5.6}
\end{equation}
which is easily evaluated using standard techniques.

Proceeding in this way, one obtains
\begin{align}
   \left.\Gamma^{(1)}_{LA}(t,p)_{\mu\nu}^{ab}
   \right|_{\rm pole}
   =g_0^2\delta^{ab}\delta(t)\delta_{\mu\nu}
   {N\over16\pi^2\eps}
   \left({3\over4}+{3\over 4\lambda_0}\right),
   \label{eq:5.7}\\[3.5pt]
   \left.\Gamma^{(1)}_{LB}(t,s,p)_{\mu\nu}^{ab}
   \right|_{\rm pole}
   =g_0^2\delta^{ab}\delta(t)\delta(s)\delta_{\mu\nu}
   {N\over16\pi^2\eps}\left(-{1\over 2\lambda_0}\right),
   \label{eq:5.8}
\end{align}
for the divergent parts of the $LA$ and $LB$ vertex functions.
The diagram $7$
and thus the $LL$ vertex function are finite at $D=4$.

\subsection{Renormalization}

The bare coupling and gauge-fixing parameter are
related to the renormalized parameters $g$ and $\lambda$
through
\begin{equation}
   g_0^2=\mu^{2\eps}g^2Z,
   \qquad
   \lambda_0=\lambda Z_3^{-1},
   \label{eq:5.9}
\end{equation}
where $\mu$ is the normalization mass.
To one-loop order, the renormalization constants
$Z$ and $Z_3$ are given by
\begin{align}
   Z=1-{b_0\over\eps}g^2+\rmO(g^4),
   \qquad
   b_0={N\over16\pi^2}{11\over3},
   \label{eq:5.10}\\[3pt]
   Z_3=1+{c_0\over\eps}g^2+\rmO(g^4),
   \qquad
   c_0={N\over16\pi^2}\left({13\over6}-{1\over2\lambda}\right),
   \label{eq:5.11}
\end{align}
up to scheme-dependent finite terms.

The renormalization of the fundamental gauge field,
\begin{equation}
   A_{\mu}^a=Z^{1/2}Z_3^{1/2}(\Ar)^a_{\mu},
   \label{eq:5.12}
\end{equation}
involves both renormalization constants
as a result of the unconventional normalization conventions
adopted in this paper. For reasons explained in sect.~\ref{sec:fin},
the ghost fields are renormalized asymmetrically according to
\begin{equation}
   c^a=\tilde{Z}_3Z^{1/2}Z_3^{1/2}(\ren{c})^a,
   \qquad
   \cbar^a=Z^{1/2}Z_3^{-1/2}(\ren{\cbar})^a,
   \label{eq:5.13}
\end{equation}
where $\tilde{Z}_3$ is the usual ghost renormalization constant.
Note that eq.~\eqref{eq:5.13} is equivalent
to the standard renormalization prescription for the
correlation functions at flow time zero, because the fields
$c$ and $\cbar$ always occur in pairs and only the
product of their renormalization factors matters.

\subsection{Do the bulk fields require renormalization?}

To one-loop order of perturbation
theory, the question may be answered by explicitly calculating
the divergent parts (if any) of bulk-field correlation functions.
We first consider the
two-point function of the $B$ field and
define its Lorentz-invariant parts $\mathcal{A}$ and $\mathcal{B}$
through
\begin{align}
   \bigl\langle\tilde{B}_{\mu}^a(t,p)\tilde{B}_{\nu}^b(s,q)\bigr\rangle
   =(2\pi)^D\delta(p+q){\delta^{ab}\over(p^2)^2}
   \notag\\[2.5ex]
   \times
   \bigl\{(\delta_{\mu\nu}p^2-p_{\mu}p_{\nu})\mathcal{A}(t,s,p^2)
   +p_{\mu}p_{\nu}\mathcal{B}(t,s,p^2)\bigr\}.
   \label{eq:5.14}
\end{align}
All self-energy diagrams drawn in fig.~\ref{fig:oneloop} contribute to
$\mathcal{A}$ and $\mathcal{B}$.
The diagrams $4-6$ actually make
two contributions, because their external legs are different.

In the renormalized perturbation expansion
\begin{equation}
   \mathcal{X}=\mu^{2\eps}\sum_{l=0}^{\infty}g^{2l+2}\mathcal{X}^{(l)},
   \qquad
   \mathcal{X}=\mathcal{A}\;\hbox{or}\;\mathcal{B},
   \label{eq:5.15}
\end{equation}
the leading-order coefficients are
\begin{equation}
   \mathcal{A}^{(0)}=\rme^{-(t+s)p^2},
   \qquad
   \mathcal{B}^{(0)}=\lambda^{-1}\rme^{-\alpha_0(t+s)p^2}.
   \label{eq:5.16}
\end{equation}
At the next order, the residues of the poles that
derive from the renormalization of the coupling and the gauge-fixing
parameter are thus given by
\begin{align}
   \left.\mathrm{res}\{\mathcal{A}^{(1)}\}\right|_{Z\;\mathrm{factors}}=
   -b_0\mathcal{A}^{(0)},
   \label{eq:5.17}\\[2.5pt]
   \left.\mathrm{res}\{\mathcal{B}^{(1)}\}\right|_{Z\;\mathrm{factors}}=
   (c_0-b_0)\mathcal{B}^{(0)}.
   \label{eq:5.18}
\end{align}
Recalling the results obtained in subsect.~\ref{ssec:selfenergy},
it is now straightforward to verify that these poles
are canceled by the divergent parts of the one-loop diagrams.
Up to this order of perturbation theory,
the two-point function of the time-dependent gauge field
is thus finite and does not require further renormalization.

In the case of the $BL$, $B\ren{A}$, $L\ren{A}$ and $LL$ correlation
functions, the cancellation of the singularities at $D=4$ can be
shown in the same way. There is actually little to prove in
the last two instances, because these two-point functions vanish
to all orders of perturbation theory (there are no diagrams
with ingoing and no outgoing flow lines). All calculations
reported in this section thus support the conjecture that
the $B$ and the $L$ field do not need to be renormalized.
